
Elevator group control in high-rise and mixed-use buildings must maintain acceptable service quality under demand that is strongly
time-varying and regime-dependent (e.g., morning up-peak, lunch inter-floor, and evening down-peak).%
\cite{cibse2020guided,strakosch2010vertical,barney2015elevatortraffic}
When mechanical capacity is fixed, a practical lever is \emph{software-level} supervision: in particular, proactively repositioning
\emph{idle} cars can reduce response distance and stabilize waiting times during regime shifts.

This report designs a deployable supervisory policy using \emph{only} the provided operational logs. We build a closed-loop pipeline that
(i) forecasts short-horizon hall-call demand, (ii) identifies the current traffic mode from real-time features, and (iii) applies
mode-aware dynamic parking for idle cars. The core parking decision is posed as a one-dimensional, demand-weighted facility-location problem
(1D $k$-median), yielding parking floors that minimize expected response distance to near-future calls. The final deliverables include both
a technical solution and a management memo describing implementation and operational safeguards.

To balance performance with deployability, we implement two complementary routes:
\textbf{Route A} (data-driven) combines learned forecasting and unsupervised mode discovery with optimization-based parking, targeting best
service performance; \textbf{Route B} (interpretable) provides a transparent baseline and a conservative fallback policy for benchmarking and
deployment risk control. Throughout, our policy is \emph{non-intrusive}: it acts only on idle cars and can sit on top of standard dispatch rules.
In particular, Route B instantiates a rule-based mode classifier and a state-aware zone parking allocator (Lobby/Mid/Upper) that generates a minimal reposition plan under simple rate limits, making it suitable both as a benchmark and as a rollback controller.

Our main contributions are:
\begin{itemize}
  \item A reproducible modeling pipeline that maps logs to a real-time traffic representation and a mode-aware parking decision.
  \item A lightweight 1D $k$-median parking formulation that is interpretable, fast, and compatible with supervisory deployment.
  \item A dual-route design (performance + fallback) and an evaluation protocol that emphasizes peak periods and regime transitions where
        proactive parking matters most.
\end{itemize}

The remainder of the paper is organized as follows. Section~\ref{sec:problem} formalizes the tasks and assumptions.
Section~\ref{sec:data} describes data processing and feature construction. Section~\ref{sec:modeling} presents the forecasting-mode
identification-parking pipeline. Section~\ref{sec:results} reports performance and robustness analyses, and Section~\ref{sec:memo}
summarizes deployment guidance in a management memo.
